\scanpage{117}
\mspage{123}
donc aussi celle de $\unsure{\mathbb D|D}$-équiv. homologique. Donc dans ce cas

\diagram{\begin{tikzcd}[ampersand replacement=\&,column sep=large]
\mathbb D\text{-fibration coh.}\arrow[r,Leftrightarrow]\arrow[d,Rightarrow] \& \mathbb D\text{-fibration }1\text{-complète}\\
\mathbb D\text{-fibration hom.} \& 
\end{tikzcd}},

et il n'y a pas à chercher plus loin que la $\mathbb D$-fibration coh.

On pourrait songer, il est vrai, à définir \deleted{\unsure{même dans ce cas}} \inserted{(dans le cas général)} une notion plus forte, en exigeant que pour tout $p:Y''\unsure{\mathord{-}|\to}Y'$ \inserted{sur $Y$} qui soit une $\mathbb D$-équiv. $1$-complète, $q:X''-X'$ est aussi une $\mathbb D$-équiv. $1$-complète. Cette fois-ci, la notion est stable par composition, mais elle me paraît néanmoins inutilisable (en dehors du cas qu'on vient de dire, où elle coïncide avec la notion de $\mathbb D$-fibration coh., ou encore de $\mathbb D$-fibration forte), car n'étant pas en mesure de décrire les $\mathbb D$-équiv. $1$-complètes, je suis aussi incapable de \deleted{\unsure{en décri|en décr}} donner des exemples \inserted{non triviaux} de $\mathbb D$-fibrations ``strictement $1$-complètes''.

\figure{\begin{tikzpicture}\draw (0,0) -- (.19,-.24);\end{tikzpicture}}

Pour vérifier qu'\unsure{une|un} $f:X-Y$ dans Diag est une $\mathbb D$-fibration cohomologique, il faut
% archive sheet number: 117
